% ------------------------------------------------------------------------
% file `1_E_juin-connaitre-1-exercise-body.tex'
%   in folder `xsim/'
%
%     exercise of type `connaitre' with id `1'
%
% generated by the `connaitre' environment of the
%   `xsim' package v0.21 (2022/02/12)
% from source `1_E_juin' on 2024/06/20 on line 198
% ------------------------------------------------------------------------
    Coche la seule bonne réponse.
    \begin{enumerate}
        \item 24 \hspace{1em} \ldots \hspace{1em}12
              \begin{qcm}
                  \item est un diviseur de
                  \item est un multiple de
                  \item divise
              \end{qcm}
        \item De quelle nature est \(-3\) ?
              \begin{qcm}
                  \item un nombre carré
                  \item un nombre naturel
                  \item un nombre entier
                \end{qcm}
        \item Dans un repère cartésien, comment s'appelle l'axe vertical ?
                \begin{qcm}
                    \item l'axe des abscisses
                    \item l'axe des ordonnées
                    \item l'axe des coordonnées
                \end{qcm}
        \item Par combien doit-on multiplier un nombre pour calculer \SI{7}{\%} de ce nombre ?
                \begin{qcm}
                    \item 0,7
                    \item 0,07
                    \item 0,007
                \end{qcm}
        \item Lequel de ces nombres est un nombre premier ?
              \begin{qcm}
                  \item 12
                  \item 1
                  \item 7
              \end{qcm}
        \item Par quel nombre \num{1548} n'est-il pas divisible ?
              \begin{qcm}
                  \item 9
                  \item 4
                  \item 8
              \end{qcm}
        \item Quel est l'élément caractéristique d'une translation ?
              \begin{qcm}
                  \item Le vecteur
                  \item L'axe de symétrie
                  \item Le centre de symétrie
              \end{qcm}
        \item Comment s'appelle une droite qui partage un segment en deux parts égales, perpendiculairement ?
              \begin{qcm}
                  \item Une médiane
                  \item Une médiatrice
                  \item Une bissectrice
              \end{qcm}
        \item Quel nom donne-t-on à un triangle ayant trois côtés de longueurs différentes ?
              \begin{qcm}
                  \item Isocèle
                  \item Équilatéral
                  \item Scalène
              \end{qcm}
        \item \(\widehat{A}\) et \(\widehat{B}\) sont deux angles supplémentaires. Quelle égalité traduit cette proposition ?
              \begin{qcm}
                  \item \(\abs{\widehat{A}} + \abs{\widehat{B}} = \ang{180}\)
                  \item \(\abs{\widehat{A}} +  \abs{\widehat{B}} = \ang{90}\)
                  \item \(\abs{\widehat{A}} \cdot \abs{\widehat{B}} = \ang{180}\)
              \end{qcm}
        \item Que vaut le produit de 3 par \(-2\) ?
              \begin{qcm}
                  \item 6
                  \item 1
                  \item \(-6\)
              \end{qcm}
    \end{enumerate}
